\section{Rank-three bases and their iterations}
For a partial linear space S, write Γ(S) for Construction3, with objects
original points P, original lines L and flags (x,M), and incidence \[
P\sim L\iff P\notin L,\quad
P\sim(x,M)\iff P\in M\setminus\{x\},\quad
L\sim(x,M)\iff L\cap M=\{x\}.
\tag{B1}
\] A chamber is therefore exactly the ordered configuration \[
P-M-x-L,\qquad P\ne x,\quad M\ne L,
\tag{B2}
\] where the three successive incidences are incidences in S. Partial
linearity makes P∉L automatic. The source calls these ``paths of length
four''; (B2) specifies unambiguously the four vertices and three edges.

\subsection{Linear spaces and the necessary
classification}\label{linear-spaces-and-the-necessary-classification}

\textbf{Proposition B.Necessity (L092 Theorem1).} Let S be a thick
connected partial linear space. If Γ(S) is a residually connected
geometry and G≤Aut(S) is chamber-transitive on Γ(S), then G is
transitive on the configurations (B2), and S is a linear space, the dual
of a linear space, or a (4,3,4)-gon.

\textbf{Proof.} The first assertion is the chamber bijection (B2).

For any point p and any nonincident line N of S, N occurs as an object
in the point residue Γ(S)p.~It cannot be isolated: that residue is a
connected rank-two geometry and has objects of both types. Therefore
some flag (x,M) is incident to p and N. This gives a path p-M-x-N of
length three in S. Incident point-line pairs already have distance one,
so every point-line distance is at most three. It follows that any two
points, or any two lines, have distance at most four, by adding an
incidence at one endpoint.

There is a proper triangle in S. Fix a point p.~Its flag objects in
Γ(S)p are indexed by its collinear points x≠p, with the unique joining
line M=px. Thickness supplies at least two such points on different
lines through p.~Connectivity of the residue gives a path between their
flag objects. Whenever this path passes through an original line N, its
two consecutive flag objects correspond to two distinct points x,y on N,
each collinear with p, while p∉N. Thus p,x,y are a proper triangle.
Partial linearity rules out an incidence cycle of length four, and the
proper triangle gives one of length six. The gonality is exactly three.

Both point- and line-diameter are at most four. If the point-diameter is
at most three, all point-point distances, being even, are at most two:
every pair of points is collinear, so S is a linear space. If the
line-diameter is at most three, the dual assertion holds. In the
remaining case both diameters equal four, giving (4,3,4). ∎

This proof uses the necessary direction of the source's
local-connectivity discussion. It does not need to assume that
connectivity of a graph on collinear neighbors alone automatically
removes isolated anti-flag objects from a residue.

\textbf{Proposition B.LocalActions (the needed local-action clauses of
L092 Lemma3).} If S is a thick partial linear space and G is
chamber-transitive on Γ(S), then G is locally two-transitive on S.
Conversely, if G is locally two-transitive on an arbitrary thick
rank-two geometry S and Γ(S) is a geometry, then S is a partial linear
space. For a partial linear space, chamber transitivity is exactly
transitivity on (B2).

\textbf{Proof.} Given an ordered pair of distinct lines M,L through x,
choose P∈M∖\{x\} to obtain (B2). Chamber transitivity then proves that
Gx is transitive on ordered pairs of its incident lines. Given an
ordered pair of distinct points P,x on M, choose a different line L
through x; the same argument proves that GM is two-transitive on the
points of M.

For the converse, a flag object (x,M) extends to a chamber of Γ(S), so
some line L through x satisfies M∩L=\{x\}. The local two-transitivity at
x makes the stabilizer of x,M transitive on all other lines through x.
Hence every other such line has singleton intersection with M. Any two
distinct lines therefore meet in at most one point, which is partial
linearity. The last assertion is precisely the equivariant bijection
(B2). ∎

\subsection{Projective planes, including order
two}\label{projective-planes-including-order-two}

\textbf{Theorem B.ProjectivePlane.} For every prime power q≥2,
Γ(PG(2,q)) is a firm residually connected geometry satisfying (IP)₂. Its
original-point and original-line residues are respectively a dual affine
plane and an affine plane of order q; its flag residues are Kq,q. The
group PSL(3,q) is chamber-transitive.

\textbf{Proof.} This is the Steiner construction S(2,q+1,q²+q+1). Its
parameter r=q+1 satisfies the firmness criterion. Equivalently, the
three panel sizes can be read directly from (B1): q+1 completions of an
anti-flag P,L, q choices of a new line through x different from M, and q
points of M other than x. All are at least two.

At an original line L, the remaining points are the projective points
outside L. Every other projective line M meets L in a unique point x and
hence supplies exactly one flag object (x,M). The resulting incidence is
the affine plane obtained by deleting L. At an original point p, the
flag objects are indexed by all other projective points x, and the
remaining original lines avoid p.~Dualizing the projective plane gives
the corresponding dual affine plane. At a flag (x,M), the remaining
points are M∖\{x\}, the remaining lines are the q other lines through x,
and every cross pair is incident in Γ(S). These are precisely the
asserted rank-two residues, and they are connected.

The full incidence graph is connected. Two original points have an
original line containing neither, and two original lines have an
original point on neither: the complement in either count has
q²−q\textgreater0 elements. Thus the subgraph on the first two types is
connected, and every flag object has a neighbor there. This proves
residual connectivity. The affine residues are partial linear spaces and
the flag residues generalized digons, proving (IP)₂.

For transitivity choose projective coordinates sending x to ⟨e₁⟩, P to
⟨e₂⟩, M to ⟨e₁,e₂⟩ and L to ⟨e₁,e₃⟩. Every configuration (B2) admits
such a projective frame. A linear map between two frames can be made
determinant one by rescaling the third basis vector; this does not
change any of the four displayed projective objects. Therefore SL(3,q),
and its projective image PSL(3,q), acts transitively. ∎

This proves L092 Theorem3 and includes q=2, which is required for the
derived base of S(3,4,8). The source's restricted q≠2 statement alone
would not justify that application. The object counts are \[
q^2+q+1,\quad q^2+q+1,\quad(q^2+q+1)(q+1),
\] and the panel sizes q+1,q,q recover its local orders and
affine-residue diagram.

\subsection{Affine planes}\label{affine-planes}

\textbf{Theorem B.AffinePlane.} For every prime power q≥3, Γ(AG(2,q)) is
a firm residually connected geometry satisfying (IP)₂. The affine
special linear group \(\mathbb F_q^2\)⋊SL(2,q) acts regularly on its chambers. In
particular every group with the translation subgroup and linear part
containing SL(2,q) is chamber-transitive.

\textbf{Proof.} The linear space is S(2,q,q²), with r=q+1. The three
panel sizes are q,q,q−1, so the construction is firm. Its flag residues
are Kq−1,q.

The point residue at the origin is the affine plane with the origin and
every line through it removed. Its point-type objects are all nonzero
vectors, and its line-type objects are the lines avoiding the origin.
Replace each remaining line by the clique on its points. Two nonzero
vectors are adjacent in this graph exactly when they lie on different
rays through the origin. It is the complete (q+1)-partite graph with
parts of size q−1, so it is connected. Every remaining line has q
points, and every point belongs to q remaining lines; the incidence
graph is therefore connected as well.

For the line residue, take L:y=0. Its point objects are (x,y) with y≠0,
and its line objects are x=ay+b, with their point on L removed. Two
points at different nonzero y-levels lie on a unique remaining line. If
they have the same y-level, choose any point at a different nonzero
level and use two such joining lines. This is possible because q−1≥2.
Hence the point objects lie in one component, and every remaining line
has q−1≥2 points, so the whole residue is connected.

The subgraph on original points and original lines, with nonincidence as
adjacency, is connected: a line can avoid any two prescribed points, and
a point can avoid any two prescribed lines, by q≥3. Every flag object
attaches to this subgraph. Together with the already connected rank-two
residues this proves residual connectivity. The two kinds of point/line
residues are subincidence structures of the affine plane and hence are
partial linear spaces; the flag residues are generalized digons.

For transitivity translate x in (B2) to zero. The nonzero vector u=P−x
determines M. Choose any vector v on the other direction L. Given a
second chamber, a linear map can send u to its prescribed target and the
line ⟨v⟩ to its prescribed target; rescale v's image to make the
determinant one. This gives chamber transitivity under \(\mathbb F_q^2\)⋊SL(2,q).

A chamber stabilizer in that group fixes zero, fixes u, and preserves a
complementary one-dimensional subspace. In the basis (u,v), its matrix
is diagonal with first entry one; determinant one forces the second
entry to be one. Thus it is trivial. The action is regular, and there
are q³(q²−1) chambers. ∎

The source's extra group at q=4 also has a complete elementary proof.

\textbf{Proposition B.SemilinearSixteen.} AΓL(1,16), acting on \(\mathbb F_{16}\) as a
two-dimensional affine space over F4, is regular on the chambers of
Γ(AG(2,4)).

\textbf{Proof.} Translate the center x of (B2) to zero and multiply by
the inverse of P−x to make P=1. Then M is the direction \(\mathbb F_4\). The other
line L has one of the four remaining directions, indexed by the
nonidentity elements of the cyclic quotient \(\mathbb F_{16}\)\emph{/F4}, of order
five. The stabilizer of 0 and 1 in AΓL(1,16) is generated by the
Frobenius a↦a². On the direction quotient it acts by multiplication by
two on exponents modulo five, a four-cycle on the nonzero exponents.
Hence it is regular on the four possible L. The whole action is
transitive, and its order 16·15·4=960 equals the chamber
count4³(4²−1)=960. ∎

These statements pay every group-action clause of L092 Theorem4.

\subsection{The two affine residue
families}\label{the-two-affine-residue-families}

Put Tq=Γ(AG(2,q))p and Sq=Γ(AG(2,q))L, with the identifications just
proved.

In Tq, the nonzero vectors on a ray form an intrinsic noncollinearity
class. For an anti-flag (P,N), the number of points of N collinear with
P is q if N is parallel to the ray 0P, and q−1 otherwise. In the latter
case N meets that ray in exactly one retained point. Both cases occur
for q≥3, so Tq is not a partial geometry.

In Sq, use coordinates with L:y=0. A line N has q−1 retained points. For
an anti-flag (P,N), exactly one of those points has P's y-coordinate,
and is not collinear with P in the retained geometry. All the others are
joined to P by nonhorizontal lines. Thus the partial-geometry parameter
is constantly q−2.

For q\textgreater3, both families have point- and line-diameter four and
gonality three, as claimed in the source. The anti-flag counts just
computed show that every point-line distance is at most three. Each
family has noncollinear point pairs and disjoint line pairs, so both
same-type diameters are exactly four. In Tq the points (1,0),(0,1),(1,1)
give a proper triangle. In Sq take (0,1),(0,t),(1,s), where 1,t,s are
three distinct nonzero field elements; all three joining lines remain,
and the triangle is proper. Partial linearity then gives gonality three.

For q=3, Sq instead is the incidence geometry of the complete bipartite
graph K3,3: its two nonzero y-levels give the two sets of three
vertices, and each remaining affine line joins one vertex from each
level. Its incidence graph is the subdivision of K3,3, with point/line
diameters four and gonality four. This accounts for the exceptional
rank-two residue in the q=3 diagram of Theorem4.

\textbf{Theorem B.PointIteration (L092 Theorem5).} For every q≥3, Γ(Tq)
is not chamber-transitive under its type-preserving automorphism group.

\textbf{Proof.} For a point P of Tq, let R(P) be P together with the
other points not collinear with it. This is exactly its nonzero ray
through the deleted origin and is defined by the incidence structure
alone.

Every type-preserving automorphism of Γ(Tq) restricts on its original
point and line types to an automorphism of Tq, because original
incidence is the complement of their incidence in (B1). Therefore it
preserves the integer \textbar N∩R(P)\textbar{} for a flag (P,N) of
Γ(Tq).

There are two possible values, zero and one, and both occur. Set
P=(1,0). The line N:b=1 avoids R(P), giving zero. For any c∈Fq*∖\{1\},
the line N:a=c does not contain P and meets R(P) at (c,0), giving one.
Such c exists because q≥3.

Both flags extend to chambers. For any anti-flag P,N, there are q net
lines through P; excluding a possible one parallel to N still leaves at
least q−1 choices meeting N. Their intersection x and the chosen line M
give the required flag object (x,M). Thus chamber transitivity would
imply transitivity on these two-object flags, contrary to the preserved
integer. ∎

This proof does not rely on the source's asserted full
automorphism-group identification for Tq. It also avoids an ambiguity
about ``paths of length four.''

The second iteration Γ(Sq) is treated in the next section. Its
exceptional A-geometry is q=8. The printed value q=4 fails residual
connectivity; a finite Lean certificate verifies a nonconstant invariant
of its literal residue graph, while the written proof below identifies
the three hexagons.

\subsection{Higher-dimensional spaces}\label{higher-dimensional-spaces}

\textbf{Theorem B.SpaceObstruction (L092 Theorems7--8, with
characteristic restrictions removed).} Let n≥3. For every prime power q,
Γ(PG(n,q)) is not residually connected. For every prime power q≥3,
Γ(AG(n,q)) is not residually connected.

\textbf{Proof.} Fix an original line L. Every point P outside L
determines a unique plane Π=⟨L,P⟩. Every flag object (x,M) in the
residue at L has M∩L=\{x\}, so M≠L and also determines a unique plane
Π=⟨L,M⟩. If P and (x,M) are incident in the residue, then P∈M∖L and the
two planes coincide. Hence every edge, and every path, preserves Π.

When n≥3 there are at least two planes through L, each with points
outside L. The residue is therefore disconnected. This already proves
both assertions.

More precisely, in the projective case each component is the affine
plane obtained by removing L from a projective plane Π. In the affine
case each component is the deleted-line affine residue Sq inside Π,
proved connected above for q≥3. The number of components in either case
is \[
\frac{q^{n-1}-1}{q-1},
\] the number of one-dimensional subspaces in the quotient direction
space. ∎

The projective assertion includes q=2 and is the negative base needed by
the binary Steiner family. The affine assertion gives the negative base
for every subline design whose derived affine dimension exceeds two.
Neither application needs to redo a higher-rank connectivity
calculation.
